\chapter{Referência Rápida}
\label{apendice:referencia}
% ── índice ──────────────────────────────────────────────────────
\index{referência rápida|textbf}
\index{sintaxe|textbf}

% ─────────────────────────────────────────────────────────────────────────────
\section{Tipos}

\begin{center}
\begin{tabular}{lll}
\toprule
\textbf{Tipo} & \textbf{Exemplo} & \textbf{Notas} \\
\midrule
\texttt{Int}       & \texttt{42}, \texttt{-7}      & 64 bits, com sinal \\
\texttt{Float}     & \texttt{3.14}, \texttt{1e-5}  & IEEE 754 \\
\texttt{Bool}      & \texttt{true}, \texttt{false} & \\
\texttt{Nil}       & \texttt{nil}                  & Ausência de valor \\
\texttt{String}    & \texttt{"texto"}              & UTF-8 \\
\texttt{FString}   & \texttt{f"val: \{x\}"}        & Interpolação \\
\texttt{List}      & \texttt{[1, 2, 3]}            & Índice base 0 \\
\texttt{Dict}      & \texttt{\{"a": 1\}}           & Chave string ou int \\
\texttt{DataFrame} & \texttt{load "f.csv" as df}   & COW, colunas Float \\
\bottomrule
\end{tabular}
\end{center}

\section{Operadores}

\begin{center}
\begin{tabular}{ll}
\toprule
\textbf{Aritmético}    & \texttt{+  -  *  /  \^{}  \%} \\
\textbf{Comparação}    & \texttt{== != > < >= <=} \\
\textbf{Lógico}        & \texttt{and  or  not} \\
\textbf{Pertencimento} & \texttt{in} \\
\textbf{Pipe}          & \texttt{|>} \\
\textbf{Atribuição}    & \texttt{=  +=  -=  *=  /=} \\
\textbf{Range}         & \texttt{a..b} (exclusivo) \quad \texttt{a..=b} (inclusivo) \\
\bottomrule
\end{tabular}
\end{center}

\section{Operadores de séries temporais}

Disponíveis dentro de \lstinline{generate df col = expr}:

\begin{center}
\begin{tabular}{ll}
\toprule
\texttt{L.x}   & Lag de ordem 1 \\
\texttt{L2.x}  & Lag de ordem 2 \\
\texttt{F.x}   & Lead de ordem 1 \\
\texttt{D.x}   & Diferença primeira \\
\bottomrule
\end{tabular}
\end{center}

Constante especial: \texttt{\_n} = número de linha (base 1).

\section{Palavras-chave}

\begin{lstlisting}[style=inline, basicstyle=\ttfamily\footnotesize]
let  const  fn  return  if  else  for  in  while  break  continue
match  try  catch  and  or  not  true  false  nil
\end{lstlisting}

% ─────────────────────────────────────────────────────────────────────────────
\section{Estatística descritiva}

\begin{center}
\begin{tabular}{lll}
\toprule
\textbf{Função} & \textbf{Resultado} & \textbf{Formas} \\
\midrule
\texttt{mean(df, x)}              & $E(X)$             & lista / df / \texttt{if =} \\
\texttt{variance(df, x)}          & $\mathrm{Var}(X)$  & lista / df / \texttt{if =} \\
\texttt{sd(df, x)}                & $\sigma(X)$        & lista / df / \texttt{if =} \\
\texttt{median(df, x)}            & mediana            & lista / df / \texttt{if =} \\
\texttt{quantile(df, x, p)}       & percentil $p$      & lista / df / \texttt{if =} \\
\texttt{cov(df, x, y)}            & $\mathrm{Cov}(X,Y)$& df / \texttt{if =} \\
\texttt{corr\_pair(df, x, y)}     & $\rho(X,Y)$        & df / \texttt{if =} \\
\texttt{summarize(df, x)}         & dict completo      & + \texttt{detail=true} \\
\texttt{correlate(df, x, y, ...)} & matriz de correlação & display \\
\texttt{codebook df}              & sumário por coluna & display \\
\texttt{tabulate df col}          & frequências        & display \\
\bottomrule
\end{tabular}
\end{center}

% ─────────────────────────────────────────────────────────────────────────────
\section{Manipulação de DataFrames}

\begin{center}
\begin{tabular}{ll}
\toprule
\textbf{Função} & \textbf{Descrição} \\
\midrule
\texttt{load "f" as df}           & Carrega arquivo (CSV, parquet) \\
\texttt{save df "f"}              & Salva arquivo \\
\texttt{count(df [, cond])}       & Número de linhas \\
\texttt{filter(df, cond)}         & Filtra linhas \\
\texttt{select(df, cols...)}      & Seleciona colunas \\
\texttt{drop(df, cols...)}        & Remove colunas \\
\texttt{generate df col = expr}   & Cria/modifica coluna \\
\texttt{replace df col = e if c}  & Modifica com condição \\
\texttt{sort(df, col [desc])}     & Ordena \\
\texttt{dropna(df [, cols...])}   & Remove NaN \\
\texttt{append(df1, df2)}         & Empilha verticalmente \\
\texttt{merge(df1, df2, on=col)}  & Junção horizontal \\
\texttt{group\_by(df, col, agg)}  & Agrega por grupo \\
\texttt{collapse(df, expr, by=c)} & Colapsa por grupo \\
\texttt{pivot\_longer(...)}       & Largo → longo \\
\texttt{pivot\_wider(...)}        & Longo → largo \\
\texttt{reshape(df, ...)}         & Reshape geral \\
\bottomrule
\end{tabular}
\end{center}

% ─────────────────────────────────────────────────────────────────────────────
\section{Geradores aleatórios}

\begin{center}
\begin{tabular}{ll}
\toprule
\textbf{Função} & \textbf{Distribuição} \\
\midrule
\texttt{rnormal()}       & $N(0,1)$ — normal padrão \\
\texttt{runiform()}      & $U(0,1)$ — uniforme \\
\texttt{rbernoulli(p)}   & Bernoulli com probabilidade $p$ fixa \\
\texttt{seed(n)}         & Fixa semente PRNG para reprodutibilidade \\
\bottomrule
\end{tabular}
\end{center}

% ─────────────────────────────────────────────────────────────────────────────
\section{Regressão linear e extensões}

\begin{center}
\begin{tabular}{ll}
\toprule
\textbf{Estimador} & \textbf{Sintaxe} \\
\midrule
OLS               & \texttt{ols(y \textasciitilde{} x1 + x2, df)} \\
IV/2SLS           & \texttt{iv(y \textasciitilde{} x, \textasciitilde{} z, df)} \\
GMM               & \texttt{gmm(y \textasciitilde{} x, \textasciitilde{} z1 + z2, df)} \\
WLS               & \texttt{wls(y \textasciitilde{} x, df, weights="w")} \\
Rolling OLS       & \texttt{rols(y \textasciitilde{} x, df, window=50)} \\
\midrule
LASSO             & \texttt{lasso(y \textasciitilde{} x1+...+xk, df, alpha=0.1)} \\
Ridge             & \texttt{ridge\_reg(y \textasciitilde{} x1+...+xk, df, alpha=0.5)} \\
Elastic Net       & \texttt{enet(y \textasciitilde{} x1+...+xk, df, alpha=0.5)} \\
\midrule
Bootstrap         & \texttt{bootstrap("ols", y \textasciitilde{} x, df, reps=500)} \\
\midrule
Regressão quantílica & \texttt{qreg(y \textasciitilde{} x, df, tau=0.5)} \\
Regressão robusta    & \texttt{rlm(y \textasciitilde{} x, df)} \\
LOWESS               & \texttt{lowess(df, y, x, frac=0.3)} \\
\midrule
SUR (Zellner)     & \texttt{sur(df, y1 \textasciitilde{} x1, y2 \textasciitilde{} x2)} \\
\bottomrule
\end{tabular}
\end{center}

{\footnotesize \textbf{Regularização:} em \texttt{lasso}, \texttt{alpha} é a força de
penalização $\lambda$ (maior = mais zeros). Em \texttt{enet}, \texttt{alpha}
é o mixing L1/(L1+L2) (0 = Ridge, 1 = LASSO).}

\section{Dados em painel}

\begin{center}
\begin{tabular}{ll}
\toprule
\textbf{Estimador} & \textbf{Sintaxe} \\
\midrule
FE (within)     & \texttt{fe(y \textasciitilde{} x, df, id=unit)} \\
RE (GLS)        & \texttt{re(y \textasciitilde{} x, df, id=unit)} \\
FE com IV       & \texttt{feiv(y \textasciitilde{} x, z \textasciitilde{} z+w, df, id="unit")} \\
Hausman         & \texttt{hausman(m\_fe, m\_re)} \\
Mundlak (CRE)   & \texttt{mundlak(df, y \textasciitilde{} x, id="unit")} \\
AB-GMM          & \texttt{ab(y \textasciitilde{} x, df, id=unit, time=t, lags=2)} \\
PCSE            & \texttt{pcse(y \textasciitilde{} x, df, id=unit, time=t)} \\
GEE             & \texttt{gee(y \textasciitilde{} x, df, id=unit)} \\
xtgls           & \texttt{xtgls(y \textasciitilde{} x, df, id=unit, time=t)} \\
Modelo misto    & \texttt{mixedlm(y \textasciitilde{} x, df, id="unit")} \\
Logit cond. (FE)& \texttt{clogit(y \textasciitilde{} x, df, group="unit")} \\
\bottomrule
\end{tabular}
\end{center}

\section{Inferência causal}

\begin{center}
\begin{tabular}{ll}
\toprule
\textbf{Método} & \textbf{Sintaxe} \\
\midrule
DiD              & \texttt{did(y \textasciitilde{} treated + post, df)} \\
RDD              & \texttt{rd(y \textasciitilde{} x, 0.0, df)} \\
PSM              & \texttt{psm(y \textasciitilde{} d + x1, df, k=1, caliper=0.05, boot=200)} \\
Controle sintético & \texttt{synth("y", treated\_id, t0, df, id="col", time="col")} \\
\bottomrule
\end{tabular}
\end{center}

{\footnotesize \textbf{PSM:} sem \texttt{caliper}, ATT pode ser tendencioso mesmo com PSM.
\textbf{synth:} \texttt{t0} = primeiro período pós-tratamento (inclusive).}

\section{Variáveis dependentes especiais}

\begin{center}
\begin{tabular}{ll}
\toprule
\textbf{Modelo} & \textbf{Sintaxe} \\
\midrule
Logit / Probit     & \texttt{logit(y \textasciitilde{} x, df)} \quad \texttt{probit(y \textasciitilde{} x, df)} \\
Efeitos marginais  & \texttt{margins(m, df)} \\
Logit multinomial  & \texttt{mlogit(y \textasciitilde{} x, df)} \\
Logit ordenado     & \texttt{ologit(y \textasciitilde{} x, df)} \\
Probit ordenado    & \texttt{oprobit(y \textasciitilde{} x, df)} \\
\midrule
Poisson            & \texttt{poisson(y \textasciitilde{} x, df)} \\
Binomial negativa  & \texttt{nbreg(y \textasciitilde{} x, df)} \\
Zero-inflacionada  & \texttt{zinb(y \textasciitilde{} x, df)} \\
\midrule
Tobit              & \texttt{tobit(y \textasciitilde{} x, df)} \\
Heckman            & \texttt{heckman(y\_out \textasciitilde{} x, s \textasciitilde{} z, df)} \\
\midrule
GLM geral          & \texttt{glm(y \textasciitilde{} x, df, family=gamma, link=log)} \\
\bottomrule
\end{tabular}
\end{center}

\begin{center}
\begin{tabular}{ll}
\toprule
\textbf{Família GLM} & \textbf{Links disponíveis} \\
\midrule
\texttt{gaussian}  & \texttt{identity}, \texttt{log}, \texttt{inverse} \\
\texttt{binomial}  & \texttt{logit}, \texttt{probit}, \texttt{cloglog} \\
\texttt{poisson}   & \texttt{log}, \texttt{identity} \\
\texttt{gamma}     & \texttt{log}, \texttt{inverse}, \texttt{identity} \\
\texttt{inverseGaussian} & \texttt{inverse\_squared} \\
\bottomrule
\end{tabular}
\end{center}

\section{Análise de sobrevivência}

\begin{center}
\begin{tabular}{ll}
\toprule
\textbf{Método} & \textbf{Sintaxe} \\
\midrule
Kaplan-Meier & \texttt{km(time, event, df)} \\
Cox PH       & \texttt{cox(time \textasciitilde{} x1 + x2, df, event=event)} \\
\bottomrule
\end{tabular}
\end{center}

{\footnotesize \textbf{Cox:} \texttt{exp(coef)} = hazard ratio.
\texttt{event} é coluna indicadora de falha (1) / censura (0).}

\section{Séries temporais}

\begin{center}
\begin{tabular}{ll}
\toprule
\textbf{Modelo} & \textbf{Sintaxe} \\
\midrule
ARMA/ARIMA       & \texttt{arima(df, y, p=1, d=0, q=1)} \\
AR($p$) via OLS  & \texttt{autoreg(df, y, lags=2, trend="c")} \\
ARDL($p$,$q$)    & \texttt{ardl(y \textasciitilde{} x1 + x2, df, lags=p, xlags=q)} \\
GARCH            & \texttt{garch(df, y, p=1, q=1)} \\
EGARCH / GJR     & \texttt{egarch(df, y, p=1, q=1)} \quad \texttt{gjrgarch(df, y, p=1, q=1)} \\
\midrule
ADF              & \texttt{adf(df, var)} \\
KPSS             & \texttt{kpss(df, var)} \\
Phillips-Perron  & \texttt{phillips\_perron(df, var)} \\
Zivot-Andrews    & \texttt{zivot(df, var)} \\
\midrule
VAR              & \texttt{var(df, y1, y2, lags=2)} \\
VECM             & \texttt{vecm(df, y1, y2, lags=1, r=1)} \\
Cointegração     & \texttt{johansen(df, y1, y2, lags=1)} \\
Granger          & \texttt{granger(df, y1, y2, lags=2)} \\
IRF (VAR reduz.) & \texttt{irf(m\_var, steps=12)} \\
SVAR Cholesky    & \texttt{svar(df, y1, y2, lags=1, id=cholesky)} \\
SVAR longo prazo & \texttt{svar(df, y1, y2, lags=1, id=longrun)} \\
IRF estrutural   & \texttt{sirf(m\_svar, steps=12)} \\
Markov Switching & \texttt{markov(df, y, k=2, p=1)} \\
\bottomrule
\end{tabular}
\end{center}

\section{Filtros e decomposição}

\begin{center}
\begin{tabular}{lll}
\toprule
\textbf{Filtro} & \textbf{Sintaxe} & \textbf{Output} \\
\midrule
Hodrick-Prescott   & \texttt{hprescott(df, y, lambda=1600)}       & \texttt{y\_trend}, \texttt{y\_cycle} \\
Baxter-King        & \texttt{baxter\_king(df, y, low=6, high=32, k=12)} & \texttt{y\_cycle} \\
Holt-Winters / ETS & \texttt{holtwinters(df, y, trend=add, seasonal=add, period=12)} & modelo ETS \\
STL                & \texttt{stl(df, y, period=12)}               & \texttt{y\_trend}, \texttt{y\_seasonal}, \texttt{y\_resid} \\
Christiano-Fitz.   & \texttt{christiano\_fitzgerald(df, y, low=6, high=32)} & \texttt{y\_cycle} \\
\bottomrule
\end{tabular}
\end{center}

{\footnotesize \textbf{Perda de borda:} HP perde pontas; BK perde $k$ observações em cada extremo;
CF e STL mantêm todos os pontos.
\textbf{Nota:} \texttt{hamilton()} é alias do Markov Switching (Hamilton 1989),
\emph{não} do filtro Hamilton 2018.
\textbf{autoreg trend:} \texttt{"c"} = constante (padrão);
\texttt{"ct"} = const + tendência; \texttt{"t"} = só tendência; \texttt{"n"} = nenhum.
\textbf{ARDL:} fórmula padrão \texttt{y \textasciitilde{} x}; \texttt{lags} = defasagens de $y$;
\texttt{xlags} = defasagens de cada regressor; inclui $x$ contemporâneo automaticamente.}

\section{Análise multivariada}

\begin{center}
\begin{tabular}{ll}
\toprule
\textbf{Método} & \textbf{Sintaxe} \\
\midrule
PCA               & \texttt{pca(df, v1, v2, ..., n=2)} \\
Análise fatorial  & \texttt{factor(df, v1, v2, ..., n=2, rotation=varimax)} \\
Correlação canôn. & \texttt{cancorr(df, xvars=["v1","v2"], yvars=["v3","v4"])} \\
MANOVA            & \texttt{manova(df, y1, y2, by="grupo")} \\
KDE               & \texttt{kde(df, var, bwidth=0.5)} \\
\bottomrule
\end{tabular}
\end{center}

\section{Testes de hipóteses}

\begin{center}
\begin{tabular}{ll}
\toprule
\textbf{Teste} & \textbf{Sintaxe} \\
\midrule
$t$ — uma amostra   & \texttt{ttest(df, var, mu=0.0)} \\
$t$ — dois grupos   & \texttt{ttest(df, var, by=grupo)} \\
$t$ — pareado       & \texttt{ttest(df, var1, var2, paired=true)} \\
ANOVA               & \texttt{anova(df, var, by=grupo)} \\
Kruskal-Wallis      & \texttt{kruskal\_wallis(df, var, by=grupo)} \\
\midrule
Normalidade (SW)    & \texttt{swilk(df, var)} \\
Normalidade (SF)    & \texttt{sfrancia(df, var)} \\
Homocedasticidade   & \texttt{vif(m\_ols)} \\
Hausman             & \texttt{hausman(m\_fe, m\_re)} \\
Teste linear        & \texttt{test(m, "x1 = x2")} \\
Combinação linear   & \texttt{lincom(m, "x1 + x2")} \\
\bottomrule
\end{tabular}
\end{center}

\section{Apresentação de resultados}

\begin{center}
\begin{tabular}{ll}
\toprule
\textbf{Função} & \textbf{Descrição} \\
\midrule
\texttt{display m}             & Exibe resultado do modelo \\
\texttt{esttab(m1, m2, ...)}   & Tabela comparativa \\
\texttt{eststo(m, "nome")}     & Armazena modelo para \texttt{esttab} \\
\texttt{predict(m, df)}        & Adiciona coluna \texttt{yhat} ao df \\
\texttt{margins(m, df)}        & Efeitos marginais médios (AME) \\
\bottomrule
\end{tabular}
\end{center}

\section{Construção de painel (fórmulas)}

Para construir painel a partir de DataFrame empilhado ($N$ unidades, $T$ períodos):

\begin{lstlisting}[style=inline]
generate df unit = (_n - 1) % N + 1
generate df t    = (_n - (_n - 1) % N - 1) / N + 1
\end{lstlisting}

Onde \texttt{N} é o número de unidades. Substitua \texttt{N} pelo valor numérico.

\section{Atalhos do REPL}

\begin{center}
\begin{tabular}{ll}
\toprule
\textbf{Atalho}  & \textbf{Ação} \\
\midrule
\texttt{exit}    & Sair do REPL \\
\texttt{Ctrl-D}  & Sair do REPL (EOF) \\
\texttt{Ctrl-C}  & Cancelar linha atual \\
\texttt{Ctrl-R}  & Buscar no histórico \\
\bottomrule
\end{tabular}
\end{center}
